% ============================================================
% config.tex — EJEMPLO LLENO (artículo Braille)
% ============================================================

\newcommand{\Titulo}{Diseño y evaluación de una red neuronal convolucional ligera para el reconocimiento óptico de Braille Grado 1 en dispositivos móviles}
\newcommand{\AutorA}{Rony Frank Churata Huania}
\newcommand{\AutorB}{Erick Jean Ramos Vera}
\newcommand{\Filiacion}{Escuela Profesional de Ingeniería de Sistemas e Informática, Universidad Nacional de Moquegua, Moquegua, Perú}

\newcommand{\PalabrasClave}{reconocimiento óptico de Braille; red neuronal convolucional ligera; detección de objetos; evaluación de modelos; educación inclusiva}
\newcommand{\Resumen}{%
En las aulas de Educación Básica Regular de Moquegua e Ilo, los estudiantes con
ceguera escriben en Braille Grado 1 con regleta y punzón sobre cartulina, un
relieve sin contraste cromático que el docente de aula no puede leer; la
retroalimentación formativa queda diferida hasta la visita del especialista. Este
proyecto diseña y evalúa una red neuronal convolucional ligera para el
reconocimiento óptico de Braille Grado 1 en español, ejecutable sin conexión en
teléfonos de gama media/baja. Se comparan cuatro configuraciones (YOLOv8n o
YOLOv11, en FP32 o INT8) mediante mAP@0.5, F1 macro, Kappa de Cohen, AUC-ROC
macro, CER, WER, latencia, memoria y tamaño, sobre validación cruzada de cinco
pliegues agrupada por página con un corpus de 40 páginas escolares (480
fotografías). El efecto pedagógico se evalúa con un cuasiexperimento de medidas
repetidas con 30 docentes de las UGEL Ilo y Mariscal Nieto. Se espera identificar
la configuración con mejor compromiso precisión--eficiencia--robustez y aportar un
corpus anotado de Braille escolar manuscrito en español peruano.
}
